\subsection{Geometrische Definition}
\label{subsec:vector_algebra_cross_product_geometric_definition}

\begin{DefinitionBox}{Kreuzprodukt}{vector_algebra_cross_product}
  \begin{minipage}[c]{0.58\linewidth}
    Das \emph{Kreuzprodukt} \EnglishTerm{cross product} zweier Vektoren
    $\vec a,\vec b\in\R^3$ ist der Vektor
    \[
      \vec c=\vec a\times\vec b
    \]
    mit den folgenden Eigenschaften:
    \begin{itemize}[leftmargin=1.4em,itemsep=0.45em]
      \item $|\vec c|=|\vec a|\,|\vec b|\sin(\varphi)$,
      \item $\vec c\perp\vec a$ und $\vec c\perp\vec b$,
      \item $\vec a,\vec b,\vec c$ bilden in dieser Reihenfolge ein
        Rechtssystem.
    \end{itemize}
  \end{minipage}\hfill
  \begin{minipage}[c]{0.38\linewidth}
    \centering
    \CachedTikzFigure{content/resources/chapters/vector_algebra/tikz/cross_product_geometry.tex}
  \end{minipage}
\end{DefinitionBox}

Das Ergebnis des Kreuzprodukts ist somit wieder ein Vektor. Das Kreuzprodukt
ist in diesem Skript ausschließlich für Vektoren im $\R^3$ definiert.

Die Orientierung des Ergebnisvektors wird mit der \emph{Rechte-Hand-Regel}
\EnglishTerm{right-hand rule} bestimmt: Zeigt der Daumen in Richtung von
$\vec a$ und der Zeigefinger in Richtung von $\vec b$, so zeigt der
Mittelfinger in Richtung von $\vec a\times\vec b$.

\noindent
\begin{minipage}[c]{0.43\linewidth}
  \centering
  \CachedTikzFigure{content/resources/chapters/vector_algebra/tikz/cross_product_unit_vectors.tex}
\end{minipage}\hfill
\begin{minipage}[c]{0.54\linewidth}
  \centering
  \CachedTikzFigure{content/resources/chapters/vector_algebra/tikz/cross_product_right_hand_rule.tex}
\end{minipage}

\begin{KeyPointBox}{Geometrische Sonderfälle}{vector_algebra_cross_product_special_cases}
  Für $\vec a,\vec b\in\R^3$ gilt:
  \begin{itemize}[leftmargin=1.5em,itemsep=0.55em]
    \item Sind $\vec a$ und $\vec b$ kollinear, so ist
      $\vec a\times\vec b=\vec 0$.
    \item Sind $\vec a$ und $\vec b$ orthogonal, so ist
      $|\vec a\times\vec b|=|\vec a|\,|\vec b|$.
  \end{itemize}
\end{KeyPointBox}

\begin{KeyPointBox}{Kreuzprodukte der kanonischen Einheitsvektoren}{vector_algebra_cross_product_unit_vectors}
  In zyklischer Reihenfolge gilt
  \[
    \vec e_1\times\vec e_2=\vec e_3,
    \qquad
    \vec e_2\times\vec e_3=\vec e_1,
    \qquad
    \vec e_3\times\vec e_1=\vec e_2.
  \]
  Bei umgekehrter Reihenfolge kehrt sich die Richtung um, beispielsweise
  \[
    \vec e_2\times\vec e_1=-\vec e_3.
  \]
  Das Vertauschen der beiden Faktoren hat also eine erkennbare Wirkung. Diese
  Beobachtung wird in den Rechenregeln verallgemeinert.
\end{KeyPointBox}
